\section{Introduction}

The dominant paradigm for equipping large language models with memory places memory as \emph{content} in a single shared context window. RAG retrieves documents into the prompt. MemGPT manages a working-memory buffer of text. Reflexion accumulates verbal reflections as prompt prefixes. System prompts define behaviour through instructions the model reads. Methods that encode information into parameters---LoRA adapters, learned soft prompts---avoid the context window but require gradient descent per adaptation and produce fixed representations that do not update from ongoing experience. In all cases, the model's memory is either explicit content competing for attention in a single context, or frozen parameters set at training time. We are missing a third option: a \emph{non-context-mediated memory channel} that persists across forward passes, adapts from experience without weight updates, and operates outside the chain-of-thought an adversary can argue within. By analogy to dual-process accounts of human memory \citep{graf1985implicit, schacter1987implicit} --- the butcher whose perceptual priors have been reshaped by decades of experience does not consult a lookup table; his eyes simply land on different features --- we want an architectural channel that shapes \emph{how} the model processes its input rather than \emph{what} it reads. We make this analogy as motivation, not as a strict mechanistic claim (\S\ref{sec:related_work}).

The current architecture forces all memory into a single inspectable channel, with three well-known failure modes. First, \emph{dilution}: as the context window fills with conversation, task content, or adversarial filler, the memory competes for attention with everything else and its influence degrades. Second, \emph{sycophancy}: the model can inspect its own memory, and an adversarial interlocutor can argue against the stated beliefs until the model abandons them. Third, \emph{cost}: every token of memory is a token of context, incurring linear attention cost at every generation step. These failures manifest across application domains: user preference tracking that dilutes over multi-turn conversation, research agents whose convictions fold under peer pressure, agentic strategy logs that crowd out task context, and codebase knowledge that scales linearly with repository size.

We address all three failure modes by introducing a new architectural primitive---the \emph{enmeshed network}---and instantiating it as HEXIS (Hidden Enmeshed eXperiential Identity States; the name is a label, the technical claim is operational), a system that implements compiled stance memory for frozen transformers. Where existing approaches force all memory into a single context window where it competes for attention, HEXIS opens a \emph{parallel context channel}: a structured cognitive schema (the Mind Tree) is processed through the same host model in a private forward pass and compiled into low-rank modulation tensors that merge with the primary context's computation. The primary context carries the conversation. The parallel context carries the disposition. The two share the same forward pass---the enmeshed network reads the host's hidden states and writes modulations at each patched layer---but the parallel context never occupies primary context positions and never competes for the primary context's attention budget.

\paragraph{Contributions.} We make three contributions:

\textbf{(1) Enmeshed Networks (\S\ref{sec:enmeshed}).} We define a new architectural primitive: a lightweight parametric module that shares a frozen host's forward pass, reading and writing intermediate representations at each layer. We present a six-axis design space (blending function, rank, modulation targets, patching pattern, compilation strategy, temporal profile), position enmeshed networks against existing adaptation methods, and introduce the Mind Tree---a typed cognitive schema that structures the enmeshed network's private subcontext for efficient compilation and zero-shot curation.

\textbf{(2) HEXIS (\S\ref{sec:hexis}).} We instantiate Level 1 (additive) enmeshment as HEXIS, implementing compiled stance memory via low-rank Q/V modulation trained with causal necessity objectives on a frozen Qwen3.5-4B-Base model. A compiled variant processes the Mind Tree in a private forward pass and produces persistent modulation tensors that shape the host's attention geometry and value extraction without occupying any context positions. The compiled signal composes with representation engineering (a fixed directional channel $d^*$) and explicit context (a curated slot of novel content selected by the enmeshed network's activation scores and the Mind Tree's structural properties).

\textbf{(3) Empirical validation (\S\ref{sec:experiments}).} We demonstrate that compiled enmeshment is dilution-immune by construction (100\% stance at 4K tokens of filler), provides sycophancy resistance (83\% $\geq$4-conviction hold rate across 5 escalating pressure levels on the 4B instruct model vs.\ 0\% for non-compiled conditions), and carries sufficient content through a rank-16 bottleneck to flip the host model's default stance on contested topics. A three-layer architecture (compiled modulation + curated slot + recursive expansion) matches full in-context beliefs at 82\% token savings. We characterise the bottleneck boundary precisely: compiled enmeshment steers parametric knowledge but cannot inject novel content---specific numbers and proper nouns do not survive the rank-16 bottleneck. The same hidden-state bridge that compiles disposition also serves agentic tool orchestration via a contrastive retrieval phi; on $\tau^3$-bench (4 domains, 3-judge audited) the teacher-loop variant lifts pass rate by $+10$pp on a 24-task balanced primary panel.

% Figure ``Four approaches to equipping LLMs with memory'' moved to Appendix~\ref{sec:arch_figure}
% to free body-page budget; the four categories (context/parameter/activation/enmeshed) are described in
% the surrounding paragraph and revisited formally in \S\ref{sec:enmeshed}.
